% ============================================================
% Section 177: 普通原子与氢原子Stark效应的区别
% ============================================================
\section{量子力学：普通原子与氢原子Stark效应的区别}
\label{sec:stark-effect-hydrogen}

\begin{modelbox}[题目：普通原子与氢原子Stark效应的区别]
通常观察到的原子中Stark效应与场强平方成正比，解释为什么如此？但对氢原子某些态Stark效应是场强的线性函数，解释为什么。用微扰计算说明氢原子基态与第一激发态的最低不为零Stark效应。已知氢原子前几个态的波函数：
\begin{align*}
\psi_{100}&\propto e^{-r/a_0},&\psi_{200}&\propto(2a_0-r)e^{-r/2a_0},\\
\psi_{21\pm1}&\propto re^{-r/2a_0}\sin\theta\,e^{\pm i\varphi},&\psi_{210}&\propto re^{-r/2a_0}\cos\theta.
\end{align*}
\end{modelbox}

\begin{tipbox}[考点快速分析]
\begin{itemize}
    \item \textbf{普通原子}：$l$ 简并已被消除，无简并微扰一级修正为零（奇宇称），二级$\propto\mathcal{E}^2$
    \item \textbf{氢原子}：库仑对称性导致 $n$ 能级内 $l$ 简并，简并微扰一级修正$\propto\mathcal{E}$
    \item \textbf{基态}：$n=1$ 无 $l$ 简并，一级为零，二级$\propto\mathcal{E}^2$
    \item \textbf{$n=2$}：四重简并，微扰矩阵非对角，能级分裂三条
\end{itemize}
\end{tipbox}

\begin{derivationbox}[标准解题过程]
\textbf{1. 定性区别}

微扰 $H'=e\mathcal{E}z=e\mathcal{E}r\cos\theta$（奇宇称）。

\textit{普通原子}：不同 $l$ 的态能量不同。$\langle nlm|H'|nlm\rangle=0$（奇宇称），一级修正为零。二级修正$\propto|\langle n'l'm|H'|nlm\rangle|^2\propto\mathcal{E}^2$。

\textit{氢原子}（$n\ge2$）：库仑势的 $SO(4)$ 对称性使同一 $n$ 内不同 $l$ 简并。简并微扰论一级修正由久期方程决定，矩阵元$\propto\mathcal{E}$，故能移$\propto\mathcal{E}$（线性Stark效应）。

\textbf{2. 氢原子基态的Stark效应}

$|100\rangle$ 非简并，一级修正 $E_{100}^{(1)}=\langle100|e\mathcal{E}z|100\rangle=0$（奇宇称）。

二级修正：
\[
E_{100}^{(2)}=\sum_{n\ge2}\frac{|\langle n10|e\mathcal{E}z|100\rangle|^2}{E_1-E_n}=-\frac{9}{4}a_0^3\mathcal{E}^2.
\]
（负号表示能量降低。）

\textbf{3. 氢原子第一激发态的Stark效应}

$n=2$ 四重简并：$|200\rangle,|210\rangle,|211\rangle,|21,-1\rangle$。

选择定则：$H'$ 为奇宇称、$\Delta m=0$，故仅 $\langle210|H'|200\rangle$ 非零。计算得
\[
\langle210|e\mathcal{E}z|200\rangle=-3ea_0\mathcal{E}\equiv-\lambda.
\]

微扰矩阵（基顺序：$|211\rangle,|21,-1\rangle,|210\rangle,|200\rangle$）：
\[
H'=\begin{pmatrix}0&0&0&0\\0&0&0&0\\0&0&0&-\lambda\\0&0&-\lambda&0\end{pmatrix}.
\]

对角化得一级修正：
\begin{equation}
\boxed{E^{(1)}=+3ea_0\mathcal{E},\;-3ea_0\mathcal{E},\;0,\;0.}
\end{equation}

对应本征态：
\begin{itemize}
    \item $+\lambda$：$(|200\rangle-|210\rangle)/\sqrt2$
    \item $-\lambda$：$(|200\rangle+|210\rangle)/\sqrt2$
    \item $0$：$|211\rangle,\;|21,-1\rangle$（横向极化，不受电场影响）
\end{itemize}

$n=2$ 分裂为三条，能移$\propto\mathcal{E}$（线性Stark效应）。
\end{derivationbox}

\begin{formulabox}[答案汇总]
\begin{center}
\begin{tabular}{ll}
\toprule
\textbf{结论} & \textbf{结果} \\
\midrule
普通原子Stark效应 & 平方效应：$\Delta E\propto\mathcal{E}^2$ \\
氢原子 $n\ge2$ & 线性效应：$\Delta E\propto\mathcal{E}$ \\
基态二级修正 & $E_{100}^{(2)}=-9a_0^3\mathcal{E}^2/4$ \\
$n=2$ 一级分裂 & $\pm3ea_0\mathcal{E},\;0,\;0$ \\
\bottomrule
\end{tabular}
\end{center}
\end{formulabox}

\begin{insightbox}[物理图像]
\begin{itemize}
    \item 氢原子的线性Stark效应是库仑势特殊对称性的直接体现。$SO(4)$ Runge-Lenz矢量守恒导致 $l$ 简并，使不同抛物线量子数的态可混合，形成沿电场方向的永久电偶极矩。
    \item 普通原子中 $l$ 简并被自旋-轨道耦合和多电子屏蔽效应消除，不存在永久电偶极矩，只能产生诱导偶极矩（$\propto\mathcal{E}$），故能移$\propto\mathcal{E}^2$。
    \item 在极强电场下，氢原子高 $n$ 态的线性Stark效应可覆盖很宽的能量范围，是研究波包动力学和经典混沌的绝佳平台。
\end{itemize}
\end{insightbox}
